\documentclass[11pt]{article}

\usepackage{amsmath}
\usepackage{amssymb}
\usepackage{geometry}
\usepackage{bm}
\usepackage{booktabs}
\usepackage{hyperref}

\geometry{margin=1in}

\title{Overview: The Radiative Transfer Equation \\ and Chandrasekhar's Formal Solution}
\author{Notes for EE 445/645 --- Physical Models in Remote Sensing}
\date{}

\begin{document}

\maketitle

\section{The Radiative Transfer Equation}

The RTE is fundamentally a \textbf{photon balance equation}. Along a ray in direction
$\bm{\Omega}$, the change in radiance $I$ per unit path length $s$ equals gains minus losses:

\begin{equation}
\frac{dI}{ds} =
\underbrace{-\kappa_e\, I}_{\text{extinction}}
+\underbrace{\kappa_s \int_{4\pi} p(\bm{\Omega}' \to \bm{\Omega})\, I(\bm{\Omega}')\, d\Omega'}_{\text{inscattering}}
+\underbrace{\kappa_a\, B(T)}_{\text{emission}}
\end{equation}

where:
\begin{itemize}
  \item $\kappa_e = \kappa_a + \kappa_s$ \quad extinction coefficient (absorption + scattering)
  \item $p(\bm{\Omega}' \to \bm{\Omega})$ \quad phase function, normalized so that $\frac{1}{4\pi}\int_{4\pi} p\, d\Omega = 1$
  \item $B(\lambda, T)$ \quad Planck blackbody spectral radiance
  \item $\omega = \kappa_s / \kappa_e$ \quad single scattering albedo ($\omega=0$: pure absorption; $\omega=1$: pure scattering)
\end{itemize}

\subsection*{Plane-Parallel Form (Myneni's Geometry)}

For a horizontally uniform medium with $z$ increasing downward and directional cosine
$\mu = \cos\theta$, the steady-state RTE becomes:

\begin{equation}
-\mu\,\frac{\partial I_\lambda(z,\bm{\Omega})}{\partial z}
+ \sigma(z,\bm{\Omega})\, I_\lambda(z,\bm{\Omega})
= \int_{4\pi} \sigma_{s\lambda}(z,\bm{\Omega}' \to \bm{\Omega})\, I_\lambda(z,\bm{\Omega}')\, d\Omega'
\end{equation}

with boundary conditions:
\begin{align}
I_\lambda(z=0,\,\bm{\Omega}) &= I_{0\lambda}\,\delta(\bm{\Omega}-\bm{\Omega}_0) + I_{d\lambda}(\bm{\Omega}), \quad \mu < 0 \quad\text{(top: sun + diffuse sky)} \\
I_\lambda(z=z_H,\,\bm{\Omega}) &= 0, \quad \mu > 0 \quad\text{(bottom: no upwelling source below canopy)}
\end{align}

\section{Chandrasekhar's Formal Solution}

\subsection{Setup}

Consider the simplified case --- a single ray with absorption and a general source $j(s)$
(which, when scattering is present, will itself contain an integral over the radiation field):

\begin{equation}
\frac{dI}{ds} + \kappa(s)\, I = j(s)
\end{equation}

This is a first-order linear ODE. Define the \textbf{optical depth} between two points:

\begin{equation}
\tau(s_1, s_2) = \int_{s_1}^{s_2} \kappa(s')\, ds'
\end{equation}

$\tau$ is dimensionless and measures how ``opaque'' the medium is between $s_1$ and $s_2$.
A value $\tau \gg 1$ means the medium is optically thick (most radiation absorbed or scattered);
$\tau \ll 1$ means optically thin (radiation passes through relatively unimpeded).

\subsection{Derivation via Integrating Factor}

The integrating factor is:

\begin{equation}
\mu(s) = e^{\,\tau(0,\,s)} = \exp\!\left(\int_0^s \kappa(s')\, ds'\right)
\end{equation}

Multiplying both sides of the ODE:

\begin{equation}
\frac{d}{ds}\!\left[\mu(s)\, I(s)\right] = \mu(s)\, j(s)
\end{equation}

Integrating from $0$ to $s$:

\begin{equation}
\mu(s)\, I(s) - I(0) = \int_0^s \mu(s')\, j(s')\, ds'
\end{equation}

Dividing through by $\mu(s) = e^{\tau(0,s)}$ yields the \textbf{formal solution}:

\begin{equation}
\boxed{
I(s) = I(0)\, e^{-\tau(0,\,s)}
+ \int_0^s j(s')\, e^{-\tau(s',\,s)}\, ds'
}
\end{equation}

\subsection{Physical Interpretation}

Each term has a clear physical meaning:

\begin{center}
\begin{tabular}{p{5cm} p{8cm}}
\toprule
\textbf{Term} & \textbf{Physical Meaning} \\
\midrule
$I(0)\, e^{-\tau(0,s)}$ &
Boundary radiance at $s=0$, attenuated by Beer's Law as it propagates to $s$.
The factor $e^{-\tau}$ is the \emph{transmittance}. \\[6pt]
$\displaystyle\int_0^s j(s')\, e^{-\tau(s',s)}\, ds'$ &
Every source element $j(s')\,ds'$ along the path contributes to $I(s)$,
weighted by how much it is attenuated ($e^{-\tau(s',s)}$) before reaching $s$. \\
\bottomrule
\end{tabular}
\end{center}

\subsection{The Scattering Complication}

When scattering is included, the source term becomes:

\begin{equation}
j(s) = \kappa_s(s) \int_{4\pi} p(\bm{\Omega}' \to \bm{\Omega})\, I(s, \bm{\Omega}')\, d\Omega'
+ \kappa_a(s)\, B(T(s))
\end{equation}

This makes the RTE \textbf{integro-differential}: the source at every point depends on the
radiation field everywhere else. The formal solution is still valid, but $j$ itself depends on
$I$, so the equation cannot be solved in closed form in general.

This is precisely why Chapter 4 of the course is devoted entirely to numerical methods:

\begin{itemize}
  \item \textbf{Successive orders of scattering} --- iterate, adding one scattering order at a time
  \item \textbf{Discrete ordinates (DOM)} --- discretize directions $\bm{\Omega}$ into a finite set; Chandrasekhar's own method
  \item \textbf{Gauss-Seidel iteration} --- solve the discretized system iteratively
  \item \textbf{Two-stream approximation} --- reduce to upward/downward fluxes only; analytically tractable
\end{itemize}

\section{Summary}

The RTE is an energy balance: changes in radiance along a ray equal inscattering and
emission gains minus extinction losses. Chandrasekhar's formal solution expresses $I(s)$ as a
superposition of the attenuated boundary condition and all source contributions along the
path, each weighted by the transmittance $e^{-\tau}$. When scattering couples all directions
together, numerical methods are required to close the system.

\end{document}
